\subsection{Local square-loss optimization}
\label{subsec:app-local-area-optimization}

\paragraph{Exact interface.}
The feasible family $\mathcal T(a,b)$ and objective $\mathcal L$ are those of
Definition~\ref{def:area-loss-family}; the public output is
Theorem~\ref{thm:local-square-loss}.

\paragraph{Variables and feasible set.}
Fix $(a,b)\in[0,1]^2$ with $\mathcal T(a,b)\ne\varnothing$.  Normalize the
distinguished vertex to $V_0$ and range over every orientation and translation
of a closed unit equilateral triangle in $\mathcal T(a,b)$.

\paragraph{Objective.}
Minimize $\mathcal L(T)$ over $T\in\mathcal T(a,b)$.

\paragraph{Exact result.}
The optimized lower bounds are \eqref{eq:min-square-loss} and
\eqref{eq:max-square-loss}.

\paragraph{Solution.}

\begin{lemma}[Local wedge reduction]
\label{lem:area-wedge}
After rotating indices to $i=0$, write
\[
 X=V_0+x(V_1-V_0)+y(V_5-V_0),
 \qquad W=\{(x,y):x\ge0,\ y\ge0\}.
\]
If a closed unit equilateral triangle $T$ contains $V_0$, then
\[
 T\cap H=T\cap W.
\]
Normalized Euclidean area is twice the corresponding $(x,y)$-coordinate
area.
\end{lemma}

\begin{proof}
The two basis vectors have length one and angle $120^\circ$, so
\[
 \lVert(x,y)\rVert^2=x^2+y^2-xy,
 \qquad
 \left\lvert\det(V_1-V_0,V_5-V_0)\right\rvert=\frac{\sqrt3}{2}.
\]
The hexagon is
\[
 H=\{0\le x,y\le2,\ |x-y|\le1\}\subset W.
\]
Conversely, if $(x,y)\in T\cap W$, then the diameter-one condition gives
$x^2+y^2-xy\le1$.  Hence
\[
 |x-y|^2\le x^2+y^2-xy\le1,
 \qquad
 \frac34x^2\le x^2+y^2-xy,
\]
and the same estimate holds for $y$.  Thus $|x-y|\le1$ and
$0\le x,y\le2/\sqrt3<2$, proving $(x,y)\in H$.  The determinant identity
proves the area assertion.
\end{proof}

\begin{lemma}[Solved local square-loss problem]
\label{lem:app-local-square-loss-calculation}
For every feasible pair $0\le a,b\le1$ and every realizing closed unit
equilateral triangle $T$,
\begin{equation}
 \mathcal L(T)\ge\min\{a,b\}^2.                  \label{eq:app-min-square-loss}
\end{equation}
If $a+b>1$, then
\begin{equation}
 \mathcal L(T)\ge\max\{a,b\}^2.                  \label{eq:app-max-square-loss}
\end{equation}
Consequently, taking the infimum over $T\in\mathcal T(a,b)$ gives the same
inequalities for $\mathcal L(a,b)$.
\end{lemma}

\begin{proof}
In the coordinates of Lemma~\ref{lem:area-wedge}, physical rotation through
$60^\circ$ sends $(x,y)$ to $(x-y,x)$.  Put
\[
 0\le t\le1,\qquad D=1-t+t^2,\qquad z=\sqrt D.
\]
After choosing one of the six directed sides, every orientation has one of
the following two forms:
\begin{equation}
 \begin{array}{c|cc}
 &\text{first side vector}&\text{second side vector}\\ \hline
 \mathrm{I}&z^{-1}(1,t)&z^{-1}(1-t,1)\\[2mm]
 \mathrm{II}&z^{-1}(t,-(1-t))&z^{-1}(1,t).
 \end{array}                                      \label{eq:area-orientation-forms}
\end{equation}
Indeed these forms cover, respectively, the physical direction sectors
$[120^\circ,180^\circ]$ and $[60^\circ,120^\circ]$; their endpoints cover
the common boundary orientations.

For Type I introduce
\[
 U=\alpha+y-tx,\qquad V=\beta+x-(1-t)y.
\]
Then $T$ is the simplex $U,V\ge0$, $U+V\le z$.  Containment of
$(0,0),(0,a),(b,0)$ gives
\begin{equation}
 \alpha\ge tb,\quad \beta\ge(1-t)a,
 \quad \alpha+\beta+(1-t)b\le z,
 \quad \alpha+\beta+ta\le z.                 \label{eq:area-type-I-containment}
\end{equation}
The two inequalities defining $W$ become
\[
 (1-t)U+V\ge(1-t)\alpha+\beta,
 \qquad U+tV\ge\alpha+t\beta.
\]
For fixed $t$, lowering $\alpha$ or $\beta$ therefore enlarges $T\cap W$.
The outside loss is minimized at
$\alpha=tb$, $\beta=(1-t)a$.  At that placement the portion outside $W$
has vertices, in the $(U,V)$-plane,
\[
 (0,0),\quad(a+tb,0),\quad(tb,(1-t)a),
 \quad(0,b+(1-t)a).
\]
Since the Jacobian of $(x,y)\mapsto(U,V)$ has absolute value $D$,
\begin{equation}
 \mathcal L(T)\ge \mathcal L_{\mathrm I}(a,b;t)
 =\frac{(1-t)a^2+tb^2+2t(1-t)ab}{D}.          \label{eq:type-I-loss}
\end{equation}
If $a\le b$, then
\[
 D(\mathcal L_{\mathrm I}-a^2)
 =t\{b^2-a^2+(1-t)(a^2+2ab)\}\ge0;
\]
the reflected factorization proves $\mathcal L_{\mathrm I}\ge b^2$ when $b\le a$.
This proves \eqref{eq:app-min-square-loss} in Type I.

Suppose now that $a+b>1$ and $a\le b$.  Set
\[
 r=\frac ab,\qquad t_0=\frac{1-r^2}{1+2r}.
\]
Since $b>1/(1+r)$, feasibility in \eqref{eq:type-I-loss} and
\eqref{eq:area-type-I-containment} excludes
$0\le t<t_0$: for $0<t<t_0$,
\[
 \left(\frac{1+r(1-t)}{1+r}\right)^2-D
 =\frac{t(1+2r)(t_0-t)}{(1+r)^2}>0,
\]
so $b+(1-t)a>z$, while $t=0$ would give $a+b\le1$.  Hence $t\ge t_0$, and
\[
 D(\mathcal L_{\mathrm I}-b^2)
 =(1-t)b^2\{r^2-1+t(2r+1)\}\ge0.
\]
If $b\le a$, put $r=b/a$ and
$t_1=(r^2+2r)/(1+2r)$.  The identical reflected calculation excludes
$t>t_1$ and gives
\[
 D(\mathcal L_{\mathrm I}-a^2)
 =ta^2\{r^2+2r-t(1+2r)\}\ge0.
\]
Thus \eqref{eq:app-max-square-loss} holds in Type I.

For Type II put
\[
 U=\alpha+(1-t)x+ty,\qquad V=\beta+tx-y.
\]
Again $T$ is the simplex $U,V\ge0$, $U+V\le z$.  Its exact reaches on the
positive $y$- and $x$-axes are
\begin{equation}
 A=\beta,\qquad B=z-\alpha-\beta,\qquad A+B\le z\le1.
 \label{eq:area-type-II-reaches}
\end{equation}
In particular Type II is impossible when $a+b>1$.  Its wedge intersection
is the quadrilateral with vertices
\[
 (0,0),\quad(B,0),\quad
 \left(\frac{B+(1-t)A}{D},\frac{A+tB}{D}\right),
 \quad(0,A).
\]
Hence
\begin{equation}
 \mathcal L(T)=\mathcal L_{\mathrm{II}}(A,B;t)
 =1-\frac{(1-t)A^2+tB^2+2AB}{D}.              \label{eq:type-II-loss}
\end{equation}
Let $S=A+B$.  If $A\le B$, write $A=rS$, $B=(1-r)S$ with
$0\le r\le1/2$.  Since $S^2\le D$, and with
\[
 C=(1-t)r^2+t(1-r)^2+2r(1-r),
\]
we have $\mathcal L_{\mathrm{II}}\ge1-C$.  The exact factorization
\[
 1-C-r^2D=(1-t)(1-2r+tr^2)\ge0
\]
gives $\mathcal L_{\mathrm{II}}\ge r^2D\ge A^2$.  Reflection gives
$\mathcal L_{\mathrm{II}}\ge B^2$ when $B\le A$.  Since $A\ge a$ and $B\ge b$,
this proves \eqref{eq:app-min-square-loss}.  The two orientation types are
exhaustive, so the lemma follows.
\end{proof}

\paragraph{Body consequence.}
The two optimized lower bounds are exactly
Theorem~\ref{thm:local-square-loss}.

\subsection{T3-like area-loss optimization}
\label{subsec:app-t3-area-optimization}

\paragraph{Exact interface.}
The input is an actual T3-like triangle and selected reach lower bounds as in
Theorem~\ref{thm:t3-direct-loss}.

\paragraph{Variables and feasible set.}
Normalize to $V_0$.  The triangle ranges only over the T3-like part of the
closed unit-triangle family; in particular $(o,n)=(2,1)$.  The translation in
Proposition~\ref{prop:t3-translation} is used only as a closed-trace majorant.

\paragraph{Objective.}
Minimize the normalized loss of the actual T3-like triangle.

\paragraph{Exact result.}
The exact nonsupercritical and actual-triangle loss conclusions are
\eqref{eq:t3-area-nonsupercritical} and \eqref{eq:t3-loss}.

\paragraph{Solution.}

\begin{lemma}[Solved T3-like loss problem]
\label{lem:app-t3-direct-loss-calculation}
Let $T_i$ be an actual T3-like V triangle, let
$0\le a_i\le A_i$ and $0\le b_i\le B_i$, and let
$0\le\mu\le1/2$.  Then
\begin{equation}
 A_i+B_i\le1.                         \label{eq:app-t3-area-nonsupercritical}
\end{equation}
Moreover,
\begin{equation}
 a_i,b_i\ge\mu
 \quad\Longrightarrow\quad
 \mathcal L(T_i)\ge2\mu-4\mu^2.                 \label{eq:app-t3-loss}
\end{equation}
\end{lemma}

\begin{proof}
The first assertion is Lemma~\ref{lem:t3-nonsupercritical}.  For the loss
bound put $T=T_i$, $a=a_i$, $b=b_i$, and $m=\mu$, and apply the
closed-trace domination in
Proposition~\ref{prop:t3-translation}.  It gives a translated Type-II
triangle $\widehat T$ of the same area with
\[
 T\cap H\subset\widehat T\cap H,
 \qquad \mathcal L(T)\ge \mathcal L(\widehat T).
\]
In the notation of that proposition, $0<t<1$,
$D=1-t+t^2$, $z=\sqrt D$, and after reflection the translated majorant has
exact reaches
\[
 A=\beta,\qquad B=z-\beta
\]
and unique wedge vertex
\[
 Q=\left(\frac{B+(1-t)A}{D},\frac{A+tB}{D}\right),
 \qquad Q_x>1,\quad Q_y\le1.
\]
Thus $z-tA>D$, so
\begin{equation}
 A<L:=\frac{z-D}{t}=\frac{z(1-z)}t.            \label{eq:t3-A-upper}
\end{equation}
Define the exact normalized loss of this majorant by
\[
 \mathcal L_{\rm maj}(z,A)
 :=\frac{tA^2+(1-t)(z-A)^2}{D}
 =\mathcal L(\widehat T).
\]
This decreases for $A<L<(1-t)z$, and therefore
\begin{equation}
 \mathcal L(T)\ge \mathcal L(\widehat T)>
 \mathcal L_{\rm maj}(z,L).                  \label{eq:t3-limit-loss}
\end{equation}

Set $r=(1-z)/t$.  Then $0<r<1/2$, and eliminating $t,z$ gives
\[
 t=\frac{1-2r}{1-r^2},\qquad
 z=\frac{r^2-r+1}{1-r^2},
\]
\[
 L=\alpha_{\rm lim}(r):=\frac{r(r^2-r+1)}{1-r^2},\qquad
 z-L=\beta_{\rm lim}(r):=\frac{r^2-r+1}{1+r}.
\]
The limiting normalized inside area and loss are
\[
 F_0(r)=\frac{(r^2-r+1)^2}{1-r^2},
 \qquad \mathcal L_0(r)=1-F_0(r).
\]
Since the exact reach $A$ is at least the required reach lower bound $a\ge m$,
\eqref{eq:t3-A-upper}
gives $\alpha_{\rm lim}(r)>m$.  Direct simplification yields
\[
 \mathcal L_0(r)-\{2\alpha_{\rm lim}(r)-4\alpha_{\rm lim}(r)^2\}
 =\frac{r^2(5r^4-8r^3+13r^2-8r+2)}{(1-r^2)^2}\ge0,
\]
because the last polynomial is at least $9r^2-8r+2\ge2/9$ on
$[0,1/2]$.  Also $\alpha_{\rm lim}(r)\le r$, and
\[
 \mathcal L_0(r)-\frac14
 =\frac{(1-2r)(2r^3-3r^2+6r-1)}{4(1-r^2)}\ge0
 \quad\text{when }\alpha_{\rm lim}(r)\ge\frac14;
\]
the cubic is increasing and has value $11/32$ at $r=1/4$.

Finally $\psi(u)=2u-4u^2$ is increasing on $[0,1/4]$ and never exceeds
$1/4$ on $[0,1/2]$.  If $\alpha_{\rm lim}(r)\le1/4$, then
$\mathcal L_0(r)\ge\psi(\alpha_{\rm lim}(r))\ge\psi(m)$; if $\alpha_{\rm lim}(r)\ge1/4$, then
$\mathcal L_0(r)\ge1/4\ge\psi(m)$.  Together with \eqref{eq:t3-limit-loss}, this proves
\eqref{eq:app-t3-loss}.  The reflected crossing is identical.
\end{proof}

\paragraph{Body consequence.}
This solved restricted optimization supplies the actual-triangle bound in
Theorem~\ref{thm:t3-direct-loss}; it makes no assertion about unrestricted
$\mathcal L(a,b)$.

\subsection{Cyclic area-loss optimization}
\label{subsec:app-cyclic-area-optimization}

\paragraph{Exact interface.}
Use the strict handoff selector of Proposition~\ref{prop:strict-handoffs} and
the two hypotheses stated in Lemma~\ref{lem:cyclic-area-loss}.

\paragraph{Variables and feasible set.}
The variables are the cyclic selector tuple
$(\xi_0,\ldots,\xi_5)$ in the product of the six open selector intervals.
The local losses range over the six actual V triangles realizing the selected
lower bounds.

\paragraph{Objective.}
Minimize $\sum_{i=0}^5\mathcal L(T_i)$ on each of the two selected-ascent
classes.

\paragraph{Exact result.}
On each of the two feasible selector classes in
Lemma~\ref{lem:cyclic-area-loss}, the exact conclusion is
\[
 \sum_{i=0}^5\mathcal L(T_i)>1.
\]

\paragraph{Solution.}

\begin{lemma}[Solved cyclic area-loss problem]
\label{lem:app-area-cyclic-calculation}
Under either exact hypothesis in Lemma~\ref{lem:cyclic-area-loss},
\[
 \sum_{i=0}^5\mathcal L(T_i)>1.
\]
\end{lemma}

\begin{proof}
Put $m=\min_i \xi_i$ and $M=\max_i \xi_i$.  The reflection
$y_i=1-\xi_{-i-1}$ swaps each V triangle's two coordinates and preserves both
alternatives and the total loss.  Hence assume $m\le1-M$.  Every selected
V triangle coordinate is then at least $m$, so
\eqref{eq:app-min-square-loss} gives loss at least $m^2$ for every V triangle.

In case (1), choose an ascent out of a minimum plateau,
$\xi_{p-1}=m<\xi_p$.  Its V triangle contains the coordinate $1-m$ and loses at least
$(1-m)^2$ by \eqref{eq:app-max-square-loss}.  A second ascent has one coordinate
strictly larger than $1/2$ and hence loses strictly more than $1/4$.  The
other four V triangles lose at least $m^2$ each, so the total loss is strictly
larger than
\[
 4m^2+(1-m)^2+\frac14
 =5\left(m-\frac15\right)^2+\frac{21}{20}>1.
\]

In case (2), rotate so that the unique ascent leaves the minimum:
$\xi_5=m<\xi_0=M$.  The supercritical V triangle loses at least $(1-m)^2$.
Lemma~\ref{lem:t3-nonsupercritical} shows that a $\Tlike$ V triangle is different
from it, and Theorem~\ref{thm:t3-direct-loss} gives that V triangle loss at least
$2m-4m^2$.  The remaining four V triangles lose at least $m^2$ each.  Their total
loss is at least
\[
 (1-m)^2+(2m-4m^2)+4m^2=1+m^2>1.
\]
\end{proof}

\paragraph{Body consequence.}
The two strict minima give Lemma~\ref{lem:cyclic-area-loss}; together with
the structural selector they yield Proposition~\ref{prop:area-branches} for
every center class whenever $N_{\mathrm{gap}}=0$.
